\paragraph{容量缺口反演与顶端仿射窗的有限层刚性}

\begin{theorem}[容量缺口的 Mellin--Stieltjes 反演完备性]\label{thm:fold-capacity-gap-mellin-stieltjes-duality}
固定分辨率 \(m\)，并定义容量缺口
\[
\mathcal R_m(T):=2^m-\mathcal C_m^{\mathrm{cont}}(T)
=\sum_{x\in X_m}\bigl(d_m(x)-T\bigr)_+,\qquad T\ge 0.
\]
对任意复数 \(s\) 满足 \(\Re s>1\)，定义
\[
S_s(m):=\sum_{x\in X_m}d_m(x)^s.
\]
则有严格恒等式
\[
\boxed{\;
S_s(m)=s(s-1)\int_{0}^{\infty}T^{s-2}\mathcal R_m(T)\,\dd T.\;}
\]
并且对任意 \(c>1\) 成立 Mellin 反演公式
\[
\boxed{\;
\mathcal R_m(T)=\frac{1}{2\pi i}\int_{c-i\infty}^{c+i\infty}
\frac{S_s(m)}{s(s-1)}\,T^{1-s}\,\dd s,
\qquad
\mathcal C_m^{\mathrm{cont}}(T)=2^m-\mathcal R_m(T).\;}
\]
因此，连续容量曲线 \(\mathcal C_m^{\mathrm{cont}}(\cdot)\) 与半平面 \(\Re s>1\) 上的全部正矩谱 \(\{S_s(m)\}\) 互相唯一决定；等价地，它们都完整刻画纤维多重度多重集 \(\{d_m(x)\}_{x\in X_m}\)。
\end{theorem}
\begin{proof}
对任意正整数 \(d\) 与 \(\Re s>1\)，有 Beta 积分恒等式
\[
d^s=s(s-1)\int_{0}^{\infty}T^{s-2}(d-T)_+\,\dd T.
\]
事实上，
\[
\int_{0}^{\infty}T^{s-2}(d-T)_+\,\dd T
=\int_{0}^{d}T^{s-2}(d-T)\,\dd T
=d^s\int_{0}^{1}u^{s-2}(1-u)\,\dd u
=\frac{d^s}{s(s-1)}.
\]
取 \(d=d_m(x)\) 并对 \(x\in X_m\) 求和，得到
\[
S_s(m)
=s(s-1)\int_{0}^{\infty}T^{s-2}\sum_{x\in X_m}(d_m(x)-T)_+\,\dd T
=s(s-1)\int_{0}^{\infty}T^{s-2}\mathcal R_m(T)\,\dd T.
\]
由于 \(X_m\) 有限且 \(\mathcal R_m(T)=0\) 当 \(T\ge M_m\) 时成立，\(\mathcal R_m\) 是紧支撑分段线性函数，上式在 \(\Re s>1\) 上绝对收敛。
因此标准 Mellin 反演适用，得到所示反演公式。

最后，由命题 \ref{prop:fold-oracle-capacity-slope-jumps-determine-histogram}，\(\mathcal C_m^{\mathrm{cont}}(\cdot)\) 已可唯一恢复纤维直方图；
而上式与反演公式分别给出
\[
\mathcal C_m^{\mathrm{cont}}\Longrightarrow \{S_s(m)\}_{\Re s>1},
\qquad
\{S_s(m)\}_{\Re s>1}\Longrightarrow \mathcal C_m^{\mathrm{cont}},
\]
从而两者与纤维多重度多重集完全等价。证毕。
\end{proof}

\begin{theorem}[顶端仿射窗控制的有限层冻结估计]\label{thm:fold-top-affine-window-controls-finite-freezing}
固定分辨率 \(m\)，并记
\[
A_m^{\max}:=\{x\in X_m:\ d_m(x)=M_m\},
\qquad
K_m:=|A_m^{\max}|,
\]
\[
M_m^{(2)}:=\max\{d_m(x):\ d_m(x)<M_m\},
\]
其中若 \(A_m^{\max}=X_m\) 则约定 \(M_m^{(2)}:=0\)。
则有：
\begin{enumerate}
  \item 对任意 \(T\in[M_m^{(2)},M_m]\)，连续容量曲线在顶端满足精确仿射公式
  \[
  \boxed{\;
  \mathcal C_m^{\mathrm{cont}}(T)=2^m-K_m(M_m-T).\;}
  \]
  \item 对任意实数 \(a>0\)，幂和满足双边界
  \[
  \boxed{\;
  K_mM_m^a
  \le
  S_a(m)
  \le
  K_mM_m^a\left(1+\frac{|X_m|-K_m}{K_m}
  \left(\frac{M_m^{(2)}}{M_m}\right)^a\right).\;}
  \]
  \item 若定义 escort 分布
  \[
  \pi_{m,a}(x):=\frac{d_m(x)^a}{S_a(m)}
  \]
  与最大层均匀分布
  \[
  U_m^{\max}:=\mathrm{Unif}(A_m^{\max}),
  \]
  则总变差距离满足
  \[
  \boxed{\;
  \|\pi_{m,a}-U_m^{\max}\|_{\mathrm{TV}}
  \le
  \frac{|X_m|-K_m}{K_m}
  \left(\frac{M_m^{(2)}}{M_m}\right)^a.\;}
  \]
\end{enumerate}
因此，顶端仿射窗 \([M_m^{(2)},M_m]\) 的三元数据 \((M_m,K_m,M_m^{(2)})\) 已足以有限层地控制高阶矩的冻结强度与 escort 集中速度。
\end{theorem}
\begin{proof}
若 \(T\in[M_m^{(2)},M_m]\)，则对 \(x\in A_m^{\max}\) 有 \(\min(d_m(x),T)=T\)，而对 \(x\notin A_m^{\max}\) 有 \(d_m(x)\le M_m^{(2)}\le T\)，故 \(\min(d_m(x),T)=d_m(x)\)。
于是
\[
\mathcal C_m^{\mathrm{cont}}(T)
=\sum_{x\notin A_m^{\max}}d_m(x)+K_mT
=2^m-K_mM_m+K_mT,
\]
即得第一条。

第二条由幂和分解
\[
S_a(m)=K_mM_m^a+\sum_{x\notin A_m^{\max}}d_m(x)^a
\]
直接推出：下界来自舍去余项；上界来自对每个 \(x\notin A_m^{\max}\) 使用 \(d_m(x)\le M_m^{(2)}\)。

再证第三条。由于 \(d_m(x)=M_m\) 在 \(A_m^{\max}\) 上常值，\(\pi_{m,a}\) 在该集合上本身就是均匀分布，故
\[
\|\pi_{m,a}-U_m^{\max}\|_{\mathrm{TV}}
=1-\pi_{m,a}(A_m^{\max}).
\]
另一方面，
\[
\pi_{m,a}(A_m^{\max})
=\frac{K_mM_m^a}{S_a(m)}
\ge
\frac{K_mM_m^a}{K_mM_m^a+(|X_m|-K_m)(M_m^{(2)})^a},
\]
于是
\[
1-\pi_{m,a}(A_m^{\max})
\le
\frac{(|X_m|-K_m)(M_m^{(2)})^a}{K_mM_m^a}
=
\frac{|X_m|-K_m}{K_m}\left(\frac{M_m^{(2)}}{M_m}\right)^a.
\]
证毕。
\end{proof}
